\documentclass[a4paper, 12pt]{extarticle}


%%%%%%%%%%%%%%%%%%%%%%%%%%%%%%
% PACKAGES
%%%%%%%%%%%%%%%%%%%%%%%%%%%%%%

%!TEX encoding = UTF8

%%%%%%%%%%%%%%%%%%%%%%%%%%%%%%%%%
% PACKAGE IMPORTS
%%%%%%%%%%%%%%%%%%%%%%%%%%%%%%%%%


\usepackage[french]{babel}
\usepackage[T1]{fontenc}
\usepackage[utf8]{inputenc}

\usepackage[
	margin=2cm,
	right=2.2cm,
	footskip=1.5cm,
]{geometry}

%ams
\usepackage{amsmath,amsfonts,amsthm,amssymb,mathtools}

% for fancy headers
\usepackage{fancyhdr}

% for table of content stuff
\usepackage{titlesec}

\usepackage{bookmark}
\usepackage{enumitem}
\usepackage[most,many,breakable]{tcolorbox}
\usepackage{varwidth,etoolbox}
\usepackage[makeroom]{cancel}
\usepackage{xcolor}
\usepackage{multicol,array}
% algorithms
\usepackage[ruled,vlined,linesnumbered]{algorithm2e}

% delays answers
% answers are from last exercise if no label is provided
% delays exercices because of exercise bank (exercices.tex file in course folder)
%\usepackage[answerdelayed, exercisedelayed, lastexercise]{exercise}
\usepackage[answerdelayed, lastexercise]{exercise}

% for slanted sfrac
\usepackage{xfrac}
% for \uline which allows underline with line breaks
\usepackage[normalem]{ulem}

\usepackage{caption, subcaption}

% margins for points (and old (p. %page%) ).
\usepackage{marginnote}

\usepackage{makecell} % pour \thead 

%index
\usepackage[texindy]{imakeidx}
\makeindex[
	columns=2, 
	title=Index, 
	options= -L french,
	%intoc,
]
\indexsetup{level=\chapter}

%virgules
\usepackage{icomma}

% for striked out implies sign (\centernot\implies)
\usepackage{centernot}

% roman numerals for \section
\renewcommand{\thesection}{\Roman{section}} 

% tikz
\usepackage{tikz}

% package systeme pour les systèmes d'équations bien alignés
\usepackage{systeme}
% add \alpha \beta \gamma as variables
% put \\ as separator so commas don't return to line
\setsysteme{eq sep = \\, align={r,r}}
% example : 
% \systeme[sort={\alpha,\beta, \gamma}]{ \alpha + 2\beta = 4, \\ \beta = 8.} 

% for polynomial long division
\usepackage{polynom}
% example : 
% \polylongdiv[style=D]{x^3 - 1 + 2x - 2x^2}{x-1}


% date
\usepackage{advdate}


%--------------------------------------------------------------------
\makeatletter
%                            \fillwithlines


% \fillwithlines takes one argument, which is either a length or \fill
% or \stretch{number}, and it fills that much vertical space with
% horizontal lines that run the length of the current line.  That is,
% they extend from the current left margin (which depends on whether
% we're in a question, part, subpart, or subsubpart) to the right
% margin.
%
% The distance between the lines is \linefillheight, whose default value
% is set with the command
%
% \setlength\linefillheight{.25in}
%
% This value can be changed by giving a new \setlength command.
%
% The thickness of the lines is \linefillthickness, whose default value
% is set with the command
%
% \setlength\linefillthickness{.1pt}
%
% This value can be changed by giving a new \setlength command.


\newlength\linefillheight
\newlength\linefillthickness
\setlength\linefillheight{.25in}
\setlength\linefillthickness{0.1pt}

\newcommand\linefill{\leavevmode
    \leaders\hrule height \linefillthickness \hfill\kern\z@}


\def\fillwithlines#1{%
  \begingroup
  \ifhmode
    \par
  \fi
  \hrule height \z@
  \nobreak
  \setbox0=\hbox to \hsize{\hskip \@totalleftmargin
          \vrule height \linefillheight depth \z@ width \z@
          \linefill}%
  % We use \cleaders (rather than \leaders) so that a given
  % vertical space will always produce the same number of lines
  % no matter where on the page it happens to start:
  \cleaders \copy0 \vskip #1 \hbox{}%
  \endgroup
}
\makeatother
%--------------------------------------------------------------------


% tikz
\usepackage{tikz, pgfplots}
\usepackage{prerex}
% i wish external worked but idk it sucks
%\usetikzlibrary{external}
%\tikzexternalize[prefix=figures/]

% for function graph
\usetikzlibrary{positioning}
\usetikzlibrary{shapes.geometric}
\usetikzlibrary{positioning}
\usetikzlibrary{shapes.misc}
\tikzset{
dot/.style = {circle, fill=#1, minimum size=5pt,
              inner sep=0pt, outer sep=0pt},
dot/.default = black % size of the circle diameter
}
\tikzset{cross/.style={cross out, draw=black, minimum size=2*(#1-\pgflinewidth), inner sep=0pt, outer sep=0pt},
%default radius will be 1pt. 
cross/.default={1pt}}

 % for braces
\usetikzlibrary{decorations.pathreplacing}
% for hashing area
\usetikzlibrary{patterns}
% tableaux var, signe
% source https://www.sqlpac.com/fr/documents/latex-package-tkz-tab-tikz-tableaux-de-signes-et-de-variations-de-fonctions.html
\usepackage{tkz-tab}



\tikzset{
	every axis/.style = {clip=true, grid style = {opacity=.5}}
}


\pagestyle{fancy}
\fancyhead[L]{Université de Mayotte}
\fancyhead[C]{\bf Contrôle Continu}
\fancyhead[R]{\today}

%%%%%%%%%%%%%%%%%%%%%%%%%%%%%%
% SELF MADE COLORS
%%%%%%%%%%%%%%%%%%%%%%%%%%%%%%

%!TEX encoding = UTF8
%%%%%%%%%%%%%%%%%%%%%%%%%%%%%%
% SELF MADE COLORS
%%%%%%%%%%%%%%%%%%%%%%%%%%%%%%


\definecolor{myg}{RGB}{56, 140, 70}
\definecolor{myb}{RGB}{45, 111, 177}
\definecolor{myr}{RGB}{199, 68, 64}
\definecolor{mytheorembg}{HTML}{F2F2F9}
\definecolor{mytheoremfr}{HTML}{00007B}
\definecolor{mylenmabg}{HTML}{FFFAF8}
\definecolor{mylenmafr}{HTML}{983b0f}
\definecolor{mypropbg}{HTML}{f2fbfc}
\definecolor{mypropfr}{HTML}{191971}
\definecolor{myexamplebg}{HTML}{F2FBF8}
\definecolor{myexamplefr}{HTML}{88D6D1}
\definecolor{myexampleti}{HTML}{2A7F7F}
\definecolor{mydefinitbg}{HTML}{E5E5FF}
\definecolor{mydefinitfr}{HTML}{3F3FA3}
\definecolor{notesgreen}{RGB}{0,162,0}
\definecolor{myp}{RGB}{197, 92, 212}
\definecolor{mygr}{HTML}{2C3338}
\definecolor{myred}{RGB}{127,0,0}
\definecolor{myyellow}{RGB}{169,121,69}
\definecolor{myexercisebg}{HTML}{F2FBF8}
\definecolor{myexercisefg}{HTML}{88D6D1}
\definecolor{doc}{RGB}{0,60,110}

% manim colors because they're beautiful
% https://docs.manim.community/en/stable/reference/manim.utils.color.manim_colors.html

\definecolor{BLACK}{HTML}{000000}\definecolor{BLUE}{HTML}{58C4DD}\definecolor{BLUE_A}{HTML}{C7E9F1}\definecolor{BLUE_B}{HTML}{9CDCEB}\definecolor{BLUE_C}{HTML}{58C4DD}\definecolor{BLUE_D}{HTML}{29ABCA}\definecolor{BLUE_E}{HTML}{236B8E}\definecolor{DARKER_GRAY}{HTML}{222222}\definecolor{DARKER_GREY}{HTML}{222222}\definecolor{DARK_BLUE}{HTML}{236B8E}\definecolor{DARK_BROWN}{HTML}{8B4513}\definecolor{DARK_GRAY}{HTML}{444444}\definecolor{DARK_GREY}{HTML}{444444}\definecolor{GOLD}{HTML}{F0AC5F}\definecolor{GOLD_A}{HTML}{F7C797}\definecolor{GOLD_B}{HTML}{F9B775}\definecolor{GOLD_C}{HTML}{F0AC5F}\definecolor{GOLD_D}{HTML}{E1A158}\definecolor{GOLD_E}{HTML}{C78D46}\definecolor{GRAY}{HTML}{888888}\definecolor{GRAY_A}{HTML}{DDDDDD}\definecolor{GRAY_B}{HTML}{BBBBBB}\definecolor{GRAY_BROWN}{HTML}{736357}\definecolor{GRAY_C}{HTML}{888888}\definecolor{GRAY_D}{HTML}{444444}\definecolor{GRAY_E}{HTML}{222222}\definecolor{GREEN}{HTML}{83C167}\definecolor{GREEN_A}{HTML}{C9E2AE}\definecolor{GREEN_B}{HTML}{A6CF8C}\definecolor{GREEN_C}{HTML}{83C167}\definecolor{GREEN_D}{HTML}{77B05D}\definecolor{GREEN_E}{HTML}{699C52}\definecolor{GREY}{HTML}{888888}\definecolor{GREY_A}{HTML}{DDDDDD}\definecolor{GREY_B}{HTML}{BBBBBB}\definecolor{GREY_BROWN}{HTML}{736357}\definecolor{GREY_C}{HTML}{888888}\definecolor{GREY_D}{HTML}{444444}\definecolor{GREY_E}{HTML}{222222}\definecolor{LIGHTER_GRAY}{HTML}{DDDDDD}\definecolor{LIGHTER_GREY}{HTML}{DDDDDD}\definecolor{LIGHT_BROWN}{HTML}{CD853F}\definecolor{LIGHT_GRAY}{HTML}{BBBBBB}\definecolor{LIGHT_GREY}{HTML}{BBBBBB}\definecolor{LIGHT_PINK}{HTML}{DC75CD}\definecolor{LOGO_BLACK}{HTML}{343434}\definecolor{LOGO_BLUE}{HTML}{525893}\definecolor{LOGO_GREEN}{HTML}{87C2A5}\definecolor{LOGO_RED}{HTML}{E07A5F}\definecolor{LOGO_WHITE}{HTML}{ECE7E2}\definecolor{MAROON}{HTML}{C55F73}\definecolor{MAROON_A}{HTML}{ECABC1}\definecolor{MAROON_B}{HTML}{EC92AB}\definecolor{MAROON_C}{HTML}{C55F73}\definecolor{MAROON_D}{HTML}{A24D61}\definecolor{MAROON_E}{HTML}{94424F}\definecolor{ORANGE}{HTML}{FF862F}\definecolor{PINK}{HTML}{D147BD}\definecolor{PURE_BLUE}{HTML}{0000FF}\definecolor{PURE_GREEN}{HTML}{00FF00}\definecolor{PURE_RED}{HTML}{FF0000}\definecolor{PURPLE}{HTML}{9A72AC}\definecolor{PURPLE_A}{HTML}{CAA3E8}\definecolor{PURPLE_B}{HTML}{B189C6}\definecolor{PURPLE_C}{HTML}{9A72AC}\definecolor{PURPLE_D}{HTML}{715582}\definecolor{PURPLE_E}{HTML}{644172}\definecolor{RED}{HTML}{FC6255}\definecolor{RED_A}{HTML}{F7A1A3}\definecolor{RED_B}{HTML}{FF8080}\definecolor{RED_C}{HTML}{FC6255}\definecolor{RED_D}{HTML}{E65A4C}\definecolor{RED_E}{HTML}{CF5044}\definecolor{TEAL}{HTML}{5CD0B3}\definecolor{TEAL_A}{HTML}{ACEAD7}\definecolor{TEAL_B}{HTML}{76DDC0}\definecolor{TEAL_C}{HTML}{5CD0B3}\definecolor{TEAL_D}{HTML}{55C1A7}\definecolor{TEAL_E}{HTML}{49A88F}\definecolor{WHITE}{HTML}{FFFFFF}\definecolor{YELLOW}{HTML}{FFFF00}\definecolor{YELLOW_A}{HTML}{FFF1B6}\definecolor{YELLOW_B}{HTML}{FFEA94}\definecolor{YELLOW_C}{HTML}{FFFF00}\definecolor{YELLOW_D}{HTML}{F4D345}\definecolor{YELLOW_E}{HTML}{E8C11C}

%%%%%%%%%%%%%%%%%%%%%%%%%%%%
% ENVIRONMENTS
%%%%%%%%%%%%%%%%%%%%%%%%%%%%

\theoremstyle{definition}

\newtheorem{theorem}{Théorème}
\newtheorem{corollaire}[theorem]{Corollaire}
\newtheorem{lemme}[theorem]{Lemme}
\newtheorem{proposition}[theorem]{Proposition}
\newtheorem{exercice}[theorem]{Exercice}
\newtheorem{exemple}[theorem]{Exemple}
\newtheorem{definition}[theorem]{Définition}
\newtheorem*{question}{Question}
\newtheorem*{preuve}{Preuve}
\newtheorem*{remarque}{Remarque}
\newtheorem*{strategie}{Stratégie}
\newtheorem*{methode}{Méthode}
\newtheorem*{notation}{Notation}
\newtheorem*{nomenclature}{Nomenclature}
\newtheorem{axiome}[theorem]{Axiome}
\newtheorem*{heuristique}{Heuristique}
\newtheorem*{motivation}{Motivation}

\newtheorem*{definition*}{Définition}
\newtheorem*{lemme*}{Lemme}
\newtheorem*{proposition*}{Proposition}
\newtheorem*{theorem*}{Théorème}
\newtheorem*{corollaire*}{Corollaire}

%%%%%%%%%%%%%%%%%%%%%%%%%%%%
% LETTERFONTS
%%%%%%%%%%%%%%%%%%%%%%%%%%%%

%!TEX encoding = UTF8

%fonts
\usepackage{libertinus,libertinust1math}
\usepackage[T1]{fontenc}

% for bold face math mode in section titles
\AddToHook{bfseries}{\boldmath}

% for calligraphic C, D, P (important to import this after the font)
\usepackage{calrsfs}
\newcommand{\D}{\mathcal{D}}
\newcommand{\Ccal}{\mathcal{C}}
\renewcommand{\P}{\mathcal{P}}

% Schwartz
\renewcommand{\S}{\mathcal{S}} % \S est le signe paragraphe normalement

% corps
\newcommand{\C}{\mathbb{C}}
\newcommand{\R}{\mathbb{R}}
\newcommand{\Rnn}{\mathbb{R}^{2n}}
\newcommand{\Z}{\mathbb{Z}}
\newcommand{\N}{\mathbb{N}}
\newcommand{\Q}{\mathbb{Q}}
\newcommand{\E}{\mathbb{E}}
\newcommand{\DD}{\mathbb{D}}
\newcommand{\K}{\mathbb{K}}
\newcommand{\F}{\mathbb{F}}
\newcommand{\U}{\mathbb{U}}
\renewcommand{\H}{\mathbb{H}}

% order notations
\DeclareRobustCommand{\O}{%
  \text{\usefont{OMS}{cmsy}{m}{n}O}%
}

% japanese bracket
\newcommand{\japb}[1]{\langle #1 \rangle}

% arrows over partial derivatives
\newcommand{\lp}{\overleftarrow{\partial}}
\newcommand{\rp}{\overrightarrow{\partial}}

% quantization
\newcommand{\h}{\hbar}
\newcommand{\Opht}{\textrm{Op}_{\h}^{t}}
\newcommand{\Op}[2][\hbar]{\textrm{Op}_{#1}^{#2}}

% omega functions
\newcommand{\omegap}[2][\rho_0]{\omega(\partial_{#1},\partial_{#2})}
\newcommand{\omegar}[2][\rho_0]{\omega(#1,#2)}

% space before semicolon
\mathcode`\;="303B

% for \Lightning
\usepackage{marvosym}

% for \warning
% 66 or 49 idk ; depends on the computer for some reason
\newcommand{\warning}{{\fontencoding{U}\fontfamily{futs}\selectfont\char 49\relax}}

% Q(\sqrt(d)) field
\newcommand{\Qsqrt}[1]{\Q\bigl(\mspace{-3mu}\sqrt{#1}\bigr)}

% closed off square root from https://en.wikibooks.org/wiki/LaTeX/Mathematics#Roots
% simple version: does not allow cbrt{x} = sqrt[3]{x}
% New definition of square root:
% it renames \sqrt as \oldsqrt
\let\oldsqrt\sqrt
% it defines the new \sqrt in terms of the old one
\def\sqrt{\mathpalette\DHLhksqrt}
\def\DHLhksqrt#1#2{%
\setbox0=\hbox{$#1\oldsqrt{#2\,}$}\dimen0=\ht0
\advance\dimen0-0.2\ht0
\setbox2=\hbox{\vrule height\ht0 depth -\dimen0}%
{\box0\lower0.4pt\box2}}



%%%%%%%%%%%%%%%%%%%%%%%%%%%%
% MACROS
%%%%%%%%%%%%%%%%%%%%%%%%%%%%

%!TEX encoding = UTF8


%%%%%%%%%%%%%%%%%%%%%%%%%%%%
% THEOREMS
%%%%%%%%%%%%%%%%%%%%%%%%%%%%

\newcommand{\thm}[3]{\begin{theorem}[#1]\label{#3}#2\end{theorem}}
\newcommand{\cor}[3]{\begin{corollaire}[#1]\label{#3}#2\end{corollaire}}
\newcommand{\lem}[3]{\begin{lemme}[#1]\label{#3}#2\end{lemme}}
\newcommand{\prop}[3]{\begin{proposition}[#1]\label{#3}#2\end{proposition}}
\newcommand{\ex}[3]{\begin{exemple}[#1]\label{#3}#2\end{exemple}}
\newcommand{\dfn}[3]{\begin{definition}[#1]\label{#3}#2\end{definition}}
\newcommand{\qs}[2]{\begin{question}[#1]#2\end{question}}

\renewcommand\qedsymbol{$\blacksquare$}
\newcommand{\pf}[2]{\begin{preuve}[#1]#2\hfill\qedsymbol\end{preuve}}

\newcommand{\nt}[1]{\begin{remarque}#1\end{remarque}}
\newcommand{\str}[1]{\begin{strategie}#1\end{strategie}}
\newcommand{\mth}[1]{\begin{methode}#1\end{methode}}
\newcommand{\ax}[3]{\begin{axiome}[#1]\label{#3}#2\end{axiome}}
\newcommand{\notations}[1]{\begin{notation}#1 \end{notation}}
\newcommand{\nomen}[1]{\begin{nomenclature}#1 \end{nomenclature}}
\newcommand{\heur}[1]{\begin{heuristique}#1\end{heuristique}}
\newcommand{\motiv}[1]{\begin{motivation}#1\end{motivation}}


%%%%%%%%%%%%%%%%%%%%%%%%%%%%
% EXERCISES
%%%%%%%%%%%%%%%%%%%%%%%%%%%%

\newcommand{\exe}[4]{
	\begin{Exercise}[#1, label=#3]
		%\marginpar{\mbox{\scriptsize(solution p.\pageref{\ExerciseLabel-Answer})}}
		#2
	\end{Exercise}
	\begin{Answer}[ref=#3]
		#4
	\end{Answer}
}

% version 1
\newcommand{\exemulticols}[5]{
	\begin{multicols}{2}
	\begin{Exercise}[#1, label=#4]
		#2
	\end{Exercise}
	\columnbreak
		#3
	\end{multicols}
	\begin{Answer}[ref=#4]
		#5
	\end{Answer}
}

% version 2
%\newcommand{\exemulticols}[5]{
%	\begin{Exercise}[#1, label=#4]
%	\begin{multicols}{2}
%		#2
%	\columnbreak
%		#3
%	\end{multicols}
%	\end{Exercise}
%	\begin{Answer}[ref=#4]
%		#5
%	\end{Answer}
%}

\newcommand{\exeTF}[1]{}

%%%%%%%%%%%%%%%%%%%%%%%%%%%%
% OPERATORS
%%%%%%%%%%%%%%%%%%%%%%%%%%%%

% deliminators
\DeclarePairedDelimiter{\abs}{\lvert}{\rvert}
\DeclarePairedDelimiter{\norm}{\lVert}{\rVert}
\DeclarePairedDelimiter{\scal}{\langle}{\rangle}

\DeclarePairedDelimiter{\ceil}{\lceil}{\rceil}
\DeclarePairedDelimiter{\floor}{\lfloor}{\rfloor}
\DeclarePairedDelimiter{\round}{\lfloor}{\rceil}

%  - others
\DeclareMathOperator{\Lap}{\mathcal{L}}
\DeclareMathOperator{\Var}{Var} % variance
\DeclareMathOperator{\Cov}{Cov} % covariance
\DeclareMathOperator{\proj}{proj} % projection
\DeclareMathOperator{\sym}{sym} % symétrique

\DeclareMathOperator{\RE}{Re} % partie réelle
\DeclareMathOperator{\IM}{Im} % partie imaginaire

\DeclareMathOperator{\pgcd}{pgcd}
\DeclareMathOperator{\ppcm}{ppcm}

\DeclareMathOperator{\ev}{ev}

\DeclareMathOperator{\tr}{tr}
\DeclareMathOperator{\sgn}{sgn}

\DeclareMathOperator{\Frac}{Frac}

\DeclareMathOperator{\End}{End}

% Lois
\DeclareMathOperator{\Bern}{\text{Bern}}
\DeclareMathOperator{\Binom}{\text{Binom}}

% I prefer the slanted \leq
\let\oldleq\leq % save them in case they're every wanted
\let\oldgeq\geq
\renewcommand{\leq}{\leqslant}
\renewcommand{\geq}{\geqslant}


%%%%%%%%%%%%%%%%%%%%%%%%%%%%
% commands
%%%%%%%%%%%%%%%%%%%%%%%%%%%%

% tel que
\newcommand{\tqs}{\text{ tels que }}
\newcommand{\tq}{\text{ tq. }}
\newcommand{\et}{\text{ et }}
\newcommand{\ou}{\text{ ou }}
\newcommand{\pourtout}{\text{ pour tout }}
\newcommand{\sct}{\text{ sachant }}
\newcommand{\id}{\text{id}}
\newcommand{\apriori}{\emph{a priori}}


% ensemble avec bigl et bigr
\newcommand{\bigset}[1]{\bigl\{ #1 \bigr\}}
\newcommand{\Bigset}[1]{\Bigl\{ #1 \Bigr\}}
\newcommand{\bigpar}[1]{\bigl( #1 \bigr)}
\newcommand{\Bigpar}[1]{\Bigl( #1 \Bigr)}

% PLUS INFTY AND MINUS INFTY WITH NO SPACE
\newcommand{\pinfty}{{+}\infty}
\newcommand{\minfty}{{-}\infty}

% vecteur flèche
\renewcommand{\vec}[1]{\overrightarrow{#1}}

% vecteur pmatrix
\newcommand{\pvec}[2]{\begin{pmatrix} #1 \\ #2 \end{pmatrix}}

% vecteur norme
%\newcommand{\norm}[1]{\left\Vert #1 \right\Vert}

% point plan
\newcommand{\point}[3]{
	#1\left(#2 ; #3 \right)
}

% \smash avant \underline pour coller la ligne au mot
% attention : ça empêche les line breaks donc c'est pas super...
\let\oldunderline\underline
\renewcommand{\underline}[1]{\oldunderline{\smash{#1}}}

% emph + index
% for \uline which allows underline with line breaks
\usepackage[normalem]{ulem}
% i add underline just for now
\newcommand{\emphindex}[1]{\emph{\uline{#1}}\index{#1}}

% tableau croisé
\newcommand{\tableaucroise}[4]{
\begin{tabular}{|c|c|c|c|}
	\cline{2-4}
	\multicolumn{1}{c|}{} & #1 \\ \hline
	#2 \\ \hline
	#3  \\ \hline
	#4  \\ \hline
\end{tabular}
}

% emdash in mathmode
\newcommand{\dash}{\text{---}}

% python minted
%\usepackage{minted} required
\newcommand{\python}[1]{
\inputminted[
		linenos,
		gobble=0,
		breaklines=true, % otherwise it breaks for no apparent reason?
		breakafter=,,
		fontsize=\normalsize,
		numbersep=5pt,
		tabsize=4, % tab ident = 4 spaces
		fontfamily=helvetica, %important pour les signes <, >
		breaklines=false,
]{python}{#1.py}
}




% add points inside environments
% example : \begin{enumerate} \item \points{2} blabla \item \points{1} bla \end{enumerate}
\newcommand{\points}[1]{
	\leavevmode\marginpar{\makebox{[#1]}}\ignorespaces%
	\addtocounter{points}{#1}%
}

%%%%%%%%%%%%%%%%%%%%%%%%%%%%
% PACKAGE-SPECIFIC ENVIRONMENTS
%%%%%%%%%%%%%%%%%%%%%%%%%%%%

% global settings for the first level in enumerate environments
\setlist[enumerate, 1]{%
	align = left, labelwidth = \parindent%, labelsep = 0pt, leftmargin = \parindent
}


%%%%%%%%%%%%%%%%%%%%%%%%%%%%%%
% EXERCISES
%%%%%%%%%%%%%%%%%%%%%%%%%%%%%%


% counter for points (title of exercices)
\usepackage{totcount}
\newtotcounter{points}

\usepackage{ifthen}
\renewcommand{\ExerciseHeader}{
	\tikz[baseline=(R.base)]\node[draw,rectangle, thick, inner sep=2pt](R) {%
		\textbf{%
			\refAnswer{\ExerciseLabel}.%
		}%
	};\!
	%
	\ifnum\ExerciseDifficulty=0
	\else
		(\theExerciseDifficulty)
	\fi
}
\renewcommand{\DifficultyMarker}{$\star$}
\renewcommand{\AnswerHeader}{
	\tikz[baseline=(R.base)]\node[draw,rectangle, thick, inner sep=2pt](R) {\textbf{\theExercise.}};\!
}
\renewcommand{\exe}[4]{
	\begin{Exercise}[title=#1, label=#3]
		\if\relax\detokenize\expandafter{\ExerciseTitle}\relax
		%\marginpar{[Bonus]}
		\else
		\marginpar{\mbox{[\ExerciseTitle]}}
		\addtocounter{points}{\ExerciseTitle}
		\fi
		#2
	\end{Exercise}
	\begin{Answer}[ref=#3]
		#4
	\end{Answer}
}
\renewcommand{\exemulticols}[5]{
	\begin{multicols}{2}
	\begin{Exercise}[title=#1, label=#4]
		\if\relax\detokenize\expandafter{\ExerciseTitle}\relax
		%\marginnote{[Bonus]}
		\else
		\marginnote{\mbox{[\ExerciseTitle]\qquad}}
		\addtocounter{points}{\ExerciseTitle}
		\fi
		#2
	\end{Exercise}
	\columnbreak
		#3
	\end{multicols}
	\begin{Answer}[ref=#4]
		#5
	\end{Answer}
}





\SetDate[24/11/2026]
\reversemarginpar
\setlength{\marginparsep}{.5cm}

\begin{document}
\fancyhead[C]{\textbf{MG303 --- Contrôle Terminal}}

\null\vspace{-20pt}
Nom : \hrulefill ~ou Code Apprenant :  \rule{3cm}{.1pt}

\begin{itemize}[label=$\bullet$]
	\item
	Matériel autorisé : aucun (pas de polycopié, pas de calculatrice).
	\item
	Les résultats du cours peuvent être utilisés sans les redémontrer.
	%Néanmoins, il n'est pas possible de faire référence aux exercices vus en classe :
	%si un problème ressemble à un exercice déjà vu, il faut le refaire.
	\item
	L'évaluation fait \pageref{lastpage} page\ifnum\getpagerefnumber{lastpage}>1 s\fi. La somme des points est \total{points}.
	\item
	Il est toujours possible d'admettre les résultats des questions précédentes, utiles aux questions suivantes.
\end{itemize}

\hrule 

\marginpar{[pts]}



%\exe{1}{
%	Considérons le \emphindex{symbole de Pochhammer} défini par
%		\begin{align*}
%			(x)_0 = 1, && \et && (x)_n = x(x-1)(x-2)\cdots(x-n+1) \text{ pour $n\in\N^*$, }
%		\end{align*}
%	ainsi que l'opérateur de différence $\Delta$ qui associe à chaque $f\in\Z[x]$ le polynôme $f(x+1)-f(x)$.
%	
%	Montrer que, pour tout $n\in\N^*$, l'identité suivante est vraie dans $\Z[x]$.
%		\[ \Delta(x)_n = n(x)_{n-1} \]
%}{exe:Pochhammer}{
%
%}
%
%\exe{1}{
%	Soit $f, g\in\Q[X]$ deux polynômes à coefficients rationnels.
%	
%	Montrer que si $f$ divise $g$ dans $\C[X]$, alors $f$ divise $g$ dans $\Q[X]$.
%}{exe:div-C-div-Q}{
%	
%}

%\vspace{20pt}

\section*{Objectif du devoir}

Dans ce devoir, nous nous intéressons au polynôme $X^{12} - 1$ de $\Z[X]$ afin d'obtenir sa décomposition en facteurs irréductibles, qui est
	\begin{align}
	X^{12} - 1 = \left(X-1\right)\left(X+1\right)\left(X^2 + X + 1\right)\left(X^2 + 1\right)\left(X^2 - X + 1\right)\left(X^4 - X^2 +1\right). \tag{équation 1} \label{eq:1}
	\end{align}
Pour cela, il sera utile de voir $X^{12}-1$ d'abord comme un polynôme à coefficients complexes, et d'étudier les ensembles suivants.
\begin{enumerate}
	\item
	L'ensemble $\U_{12}$ des racines 12-ièmes de l'unité.
	\item
	L'ensemble $\D_{12}$ des diviseurs de 12.
		\[ \D_{12} = \bigset{ d \in \N \tq d | 12 } \]
	Tout le long du devoir, nous noterons $d|12$ pour signifier $d\in\D_{12}$.
	\item
	Pour chaque $d|12$, l'ensemble $\mu_d$ des entiers entre 1 et $d$ inclus qui sont premiers avec $d$.
		\[ \mu_d = \bigset{ 1 \leq k \leq d \tq \pgcd(k, d) = 1 } \]
\end{enumerate}

\section{Préliminaires}


\exe{5}{
	Compréhension des ensembles considérés.
	\begin{enumerate}
		\item
		Donner la définition de l'ensemble $\U_{12}$.
		\item
		Décrire explicitement tous les éléments de $\U_{12}$ en utilisant l'exponentielle complexe.
		\item
		Énumérer les éléments de $\D_{12}$.
		\item
		Énumérer les éléments de $\mu_d$ pour chaque élément $d$ de $\D_{12}$.
		\item
		Vérifier que 
			\[ \sum_{d|12} |\mu_d| = 12. \]
	\end{enumerate}
}{exe:phi12-1}{
	\begin{enumerate}
		\item
		 $\U_{12} = \bigset{ z\in\C \tq z^n = 1 }$.
		\item
		$\U_{12} = \bigset{ e^{2\pi i \frac{k}{n}} ~|~ 1 \leq k \leq n }$.
		\item
		$\D_{12} = \bigset{ 1 ; 2 ; 3; 4; 6 ; 12}$.
		\item
		\begin{enumerate}
			\item[$d=1$]
			Nous avons $\mu_1 = \{ 1 \}$.
			\item[$d=2$]
			Nous avons $\mu_2 = \{ 1 \}$.
			\item[$d=3$]
			Nous avons $\mu_3 = \{ 1 ; 2 \}$.
			\item[$d=4$]
			Nous avons $\mu_4 = \{ 1 ; 3 \}$.
			\item[$d=6$]
			Nous avons $\mu_6 = \{ 1 ; 5 \}$.
			\item[$d=12$]
			Nous avons $\mu_{12} = \{ 1 ; 5 ; 7 ; 11 \}$.
		\end{enumerate}
		\item
		Nous avons bien $\sum_{d|12} |\mu_d| = 1 + 1+ 2 + 2 + 2 + 4 = 12$.
	\end{enumerate}
}


\section{Décompositions dans $\C$}

\exe{3}{
	Décompositions polynomiales classiques.
	\begin{enumerate}
		\item
		Décomposer explicitement le polynôme $X^{12}-1$ dans $\C[X]$.
		\item
		Décomposer explicitement le polynôme $X^{11} + X^{10} + \dots + X + 1$ dans $\C[X]$.
	\end{enumerate}
}{exe:phi12-2}{
	\begin{enumerate}
		\item
		$X^{12}-1 = \prod_{k=1}^{12} \left( X - e^{2\pi i \frac{k}{12}} \right)$.
		\item
		Comme $X^{11} + X^{10} + \dots + X + 1 = \frac{X^{12}-1}{X-1}$, nous avons
			\[ X^{12}-1 = \prod_{k=1}^{11} \left( X - e^{2\pi i \frac{k}{12}} \right). \]
	\end{enumerate}
}

\newpage

\section{Partitionnement des racines de l'unité}

\exe{3}{
	Nous cherchons désormais à faire le lien entre $\U_{12}$ et les ensembles $\mu_d$ où $d$ divise $12$.
	\begin{enumerate}
		\item
		Écrire les fractions 
		\begin{align*}
			\frac1{12} && \frac2{12} && \frac3{12} && \dots && \frac{11}{12} && \frac{12}{12}
		\end{align*}
		sous forme irréductible.
		\item
		Quel lien avec les ensembles $\mu_d$ peut-on faire ?
	\end{enumerate}
}{exe:phi12-3}{
	\begin{enumerate}
		\item
		Nous obtenons
		\begin{align*}
			\frac1{12} && \frac1{6} && \frac1{4} && \frac13 && \frac5{12} && \frac12 && \frac7{12} && \frac23 && \frac34 && \frac56 && \frac{11}{12} && \frac11
		\end{align*}
		\item
		Nous remarquons que les fractions sont toutes sur un dénominateur $d$ diviseur de 12.
		De plus, pour chaque dénominateur $d|12$, les numérateurs sont exactement les éléments de $\mu_d$, c'est-à-dire les entiers $k$ entre 1 et $d$ qui sont premiers avec $d$.
	\end{enumerate}
}


\exe{3}{
	Nous allons désormais partitionner les racines 12-ièmes de l'unités en sous-ensembles associés à chaque diviseur de 12.
	
	Pour chaque $d|12$, posons le polynôme à coefficients \apriori~complexes
	\begin{align}
		\phi_d(X) = \prod_{k\in\mu_d} \left(X-e^{2\pi i \frac{k}{d}}\right). \tag{définition 2} \label{eq:2}
	\end{align}
	\begin{enumerate}
		\item
		Donner le degré de $\phi_d$ pour chaque $d\in\D_{12}$.
		\item
		Montrer que
			\[ \U_{12} = \bigcup_{d|12} \bigset{ e^{2\pi i \frac{k}{d}} ~|~ k\in\mu_d }. \]
		\item
		En conclure que
			\[ X^{12} - 1 = \prod_{d|12} \phi_d(X). \]
	\end{enumerate}
}{exe:phi12-4}{
	\begin{enumerate}
		\item
		$\deg(\phi_d) = |\mu_d|$, dont la valeur peut être lue à l'exercice \ref{exe:phi12-1}.
		\item
		Les racines 12-ièmes de l'unité s'écrivent toutes $e^{2\pi i \frac{k}{12}}$, dont la fraction $\frac{k}{12}$ se réduit en $\frac{k'}{d}$ avec $d|12$  et $k'$ premier avec $d$, comme listé à l'exercice \ref{exe:phi12-3}.
		\item
		La conclusion suit de la partition des racines selon le dénominateur de la fraction irréductible associée.
	\end{enumerate}
}

\section{Conclusion}

\exe{6}{
	Vérifions désormais que les polynômes $\phi_d$, définis par la \ref{eq:2} de l'exercice \ref{exe:phi12-4}, forment la décomposition en irréductibles de $X^{12}-1$ dans $\Z[X]$, donnée par l'\ref{eq:1} du début du devoir.

	\begin{enumerate}
		\item
		Développer et réduire chaque polynôme $\phi_d$ pour retrouver les six facteurs de la décomposition de l'\ref{eq:1} du début du devoir.
		
		\emph{Il est possible d'admettre que les $\phi_d$ sont les facteurs de l'\ref{eq:1} et de passer aux questions suivantes.}
		\item
		Justifier que chaque polynôme $\phi_d$ est polynôme primitif de $\Z[X]$.
		\item
		Montrer que chaque polynôme $\phi_d$ est irréductible sur $\Q[X]$.
		\item
		Conclure que l'\ref{eq:1} est une décomposition en facteurs irréductibles dans $\Z[X]$.
	\end{enumerate}
}{exe:phi12-5}{
	\begin{enumerate}
		\item
		Nous calculons les polynômes suivants.
		\begin{enumerate}
			\item[$d=1$]
			Nous avons $\phi_1(X) = X-1$.
			\item[$d=2$]
			Nous avons $\phi_2(X) = X-e^{2\pi i \frac12} = X+1$.
			\item[$d=3$]
			Nous avons $\phi_3(X) = (X-e^{2\pi i \frac23})(X-e^{2\pi i \frac23}) = X^2 + X + 1$.
			\item[$d=4$]
			Nous avons $\phi_4(X) = (X-e^{2\pi i \frac14})(X-e^{2\pi i \frac34}) = X^2 + 1$.
			\item[$d=6$]
			Nous avons $\phi_6(X) = (X-e^{2\pi i \frac16})(X-e^{2\pi i \frac56}) = X^2 - X + 1$.
			\item[$d=12$]
			Nous avons $\phi_12(X) = (X-e^{2\pi i \frac{1}{12}})(X-e^{2\pi i \frac{5}{12}})(X-e^{2\pi i \frac{7}{12}})(X-e^{2\pi i \frac{11}{12}}) = X^4 - X^2 + 1$.
		\end{enumerate}
		\item
		Les polynômes sont tous primitifs de $\Z[X]$ car unitaires.
		\item
		Tout polynôme de degré 1 est irréductible sur $\Q[X]$.
		
		Les polynômes de degrés 2 sont irréductibles sur $\Q[X]$ car leurs racines sont complexes.
		
		Seul $\mu_{12}(X)$ est à étudier.
		Par contradiction, si $\mu_{12}$ se scinde dans $\Q[X]$, ses facteurs sont nécessairement de degré 2 car ses racines n'appartiennent pas à $\Q$.
		De plus, pour obtenir un facteur à coefficients réel, ses racines sont nécessairement conjuguées complexes l'une de l'autre.
		Un des facteurs doit donc être
			\[ \left(X-e^{2\pi i \frac{1}{12}}\right)  \left(X-e^{-2\pi i \frac{1}{12}}\right) = X^2 - \sqrt3 X + 1, \]
		qui n'est pas un polynôme de $\Q[X]$.
		\item
		D'après le lemme de Gauss, les polynômes primitifs de $\Z[X]$ irréductible dans $\Q[X]$ sont les irréductibles de $\Z[X]$.
		Nous avons donc obtenu la décomposition de $X^{12} - 1$.
	\end{enumerate}
}

\section*{Bonus}

Le résultat suivant permet de déduire que $\phi_n$ est un polynôme à coefficient entiers relatifs pour tout $n\in\N^*$ en utilisant que
	\[ \phi_n(X) = \frac{X^n - 1}{\prod\limits_{d|n, d\neq n} \phi_d(X)}. \]

\exe{2}{
	Soit $f_0\in\Z[X]$ et $g_0\in\Z[X]$ deux polynômes primitifs.
	Montrer que si $f_0|g_0$ dans $\C[X]$, alors $f_0|g_0$ dans $\Q[X]$ et puis $f_0|g_0$ dans $\Z[X]$.
}{exe:divCdivQdivZ}{
	Nous avons $g_0 = qf_0$ avec $q\in\C[X]$.
	Notons $g_0 = q' f_0 + r$ avec $q', r \in\Q[X]$ et $\deg(r)<\deg(f_0)$ la division euclidienne de $g_0$ par $f_0$ dans $\Q[X]$.
	
	Alors, en voyant la division dans $\C[X]$, nous avons
		\[ qf_0 = q'f_0 + r, \]
	et donc $r = f_0(q-q')$.
	Or, comme $\deg(r) < \deg(f_0)$, l'égalité n'est possible que lorsque $q-q' = 0$ et donc $r=0$.
	Il suit alors que $f_0$ divise $g_0$ dans $\Q[X]$.
	
	Enfin, le lemme de Gauss permet de dire que $f_0$ divise $g_0$ dans $\Z[X]$ car $f_0$ et $g_0$ sont primitifs.
	Ceci a été vu en cours mais peut être redémontré comme suit.
	À partir de $g_0 = qf_0$ avec $q\in\Q[X]$ cette fois, nous avons $g_0 = c_q q_0 f_0$ en décomposant $q=c_q q_0$ comme produit de son contenu et d'un polynôme primitif.
	Or $q_0f_0$ est primitif d'après le lemme de Gauss, et donc, par unicité du contenu, $c_q=1$ et $g_0 = q_0f_0$ dans $\Z[X]$.
}




%%%%%%%%%%%

\label{lastpage}
\newpage
\fancyhead[C]{\textbf{Solutions}}
\shipoutAnswer

\end{document}
